
\documentclass{exam}
\usepackage{../../../mypackages}
\usepackage{../../../macros}


\tcbuselibrary{theorems,skins,hooks}

\definecolor{myindbg}{HTML}{DFF5EC}
\definecolor{myindfr}{HTML}{1F6F4A}

\newtcolorbox{Indication}{
  enhanced,
  colback=myindbg,
  colframe=myindfr,
  arc=6mm,
  rounded corners,
  boxrule=0.6mm,
  left=3mm,
  right=3mm,
  top=2mm,
  bottom=2mm
}

\title{Interrogation N°6}
\author{N. Bancel}
\date{10 Février 2026}

% Mise en forme des solutions en bleu
\SolutionEmphasis{\color{blue}}
\renewcommand{\solutiontitle}{\noindent}

\begin{document}

\dsheader{Collège Lycée Suger}{Mathématiques}{Année 2025-2026}{Classe de 6ème}
{\let\newpage\relax\maketitle}
\consignesrendusujet{1h15}{1.5}

\section*{Exercice 1 - Cours (4 points)}

\begin{Indication}
  Toute réponse approximative vaudra 0 point.
\end{Indication}

\begin{questions}
  \question[1.5] \textbf{Médiatrice} Donner la définition précise de la médiatrice d’un segment. Un vocabulaire mathématique est attendu.
  \begin{solution}
          \begin{Indication}

Ce qui est attendu : 
\begin{compactitem}
\item Penser à la perpendicularité
\item Penser au milieu du segment
\item Eviter les tournures "la droite qui forme un angle droit". "La droite fait 90° d'angle". Forme un angle droit avec quoi ?
\item Des points symétriques correctement placés : avec le quadrillage, il n'y a pas vraiment d'excuse pour un mauvais placement et des traits pas droits.
\end{compactitem}

      \end{Indication}
      La médiatrice d’un segment est \textbf{la droite qui passe par le milieu de ce segment et qui est perpendiculaire à ce segment}.
      \end{solution}
  \question[1.5] \textbf{Bissectrice} Donner la définition précise de la bissectrice d’un angle. Un vocabulaire mathématique est attendu.
      \begin{solution}
                  \begin{Indication}

Ce qui est attendu : 
\begin{compactitem}
\item On ne parle pas du milieu d'un angle  
\item On ne dit pas que l'angle est coupé en deux "parties égales" (trop vague). On dit qu'il est partagé en deux angles de même mesure.
\item On évite de dire que les angles sont égaux. Au fond, c'est leur mesure qui est égale.
\end{compactitem}

      \end{Indication}   
      La bissectrice d’un angle est \textbf{la demi-droite qui partage cet angle en deux angles de même mesure}.
    \end{solution}

    \question[1] \textbf{Cercle circonscrit} Qu'est ce que le cercle circonscrit d'un triangle ABC ? Comment le construit-on ?
  \end{questions}

\pagebreak

\section*{Exercice 2 - Symétrie axiale - Tracés (8 points)}

\begin{questions}
  \question[3] Sur la figure ci-dessous, tracer le symétrique de la figure dessinée en bas par rapport à la droite $(d)$. Il n'est pas nécessaire ici de mettre de légendes (les tracés sont à faire sur le sujet).

\fig{0.8}{01.jpg}{}

\begin{solution}
      \begin{Indication}

Ce qui est attendu : 
\begin{compactitem}
\item Une figure proprement tracée
\item Des points symétriques correctement placés : avec le quadrillage, il n'y a pas vraiment d'excuse pour un mauvais placement et des traits pas droits.
\end{compactitem}

      \end{Indication}

      

      \end{solution}


\pagebreak
  \question[5] Sur la figure ci-dessous, tracer les symétriques du triangle $ABC$ et de la droite $(d_2)$ par rapport à la droite $(d_1)$. \textrr{Vous noterez toutes les légendes nécessaires} (les tracés sont à faire sur le sujet).
\fig{1}{05.jpg}{}
\end{questions}

\begin{solution}
      \begin{Indication}

Ce qui est attendu : 
\begin{compactitem}
\item Que les angles droits soient faits à l'équerre
\item Que les traits utilisés pour arriver aux figures soient moins visibles que les figures elles-mêmes (vous pouvez appuyer moins fort ou faire des pointillés)
\item Que le cercle soit bien tracé au compas
\item Que toutes les légendes demandées soient présentes : quand on fait le symétrique d'un point, on nomme le point d'origine et on nomme son symétrique. Et on fait les codages de longueur et d'angle qui nous ont permis de tracer le symétrique.
\item \textrr{Il n'y a pas de codage de longueur à faire sur une droite. Une droite n'a pas de longueur}
\item \textrr{Une droite a besoin d'un 2ème point pour être tracée. Il fallait choisir un 2ème point au hasard sur la 1ère droite, le nommer, et construire son symétrique}
\end{compactitem}

      \end{Indication}
      


\end{solution}

\pagebreak

\section*{Exercice 3 - Angles du triangle (5 points)}

\begin{Indication}
Toute démarche, toute explication de "comment vous auriez résolu l'exercice si vous aviez eu telle ou telle information" sera valorisée
\end{Indication}

Bien analyser la figure ci-dessous, et vous en servir pour répondre aux questions suivantes.

\fig{0.8}{03.jpg}{}

\begin{questions}
  \question[1.5] \just À l’aide des informations codées sur la figure, calculer la mesure de l'angle $\widehat{BAC}$
\question[0.5] Que peut-on dire des angles à la base d'un triangle isocèle ?
\question[1.5] \just En précisant le triangle dans lequel vous faites votre raisonnement, calculer la mesure de l'angle $\widehat{BAD}$
\question[1.5] \just En déduire la mesure de l'angle $\widehat{CAD}$. Préciser la nature de cet angle.
\end{questions}

\begin{solution}
xxxx
\end{solution}

\pagebreak

\section*{Exercice 4 - Longueurs (3 points)}
\begin{Indication}
Toute démarche, toute explication de "comment vous auriez résolu l'exercice si vous aviez eu telle ou telle information" sera valorisée
\end{Indication}


\fig{0.6}{04.jpg}{}

\begin{questions}
  \question[1] \just En vous basant sur la figure ci-dessus, calculer la mesure de l'angle $\widehat{IHJ}$
  \question[1] \just Que peut-on en déduire sur la nature du triangle $HIJ$ ?
  \question[1] \just En déduire le périmètre du triangle $HIJ$. \textit{Pour rappel, le périmètre d'un polygone est la somme des longueurs des côtés de ce polygone}.
\end{questions}

\end{document}
